\chapter*{\sffamily Agradecimientos}
\addcontentsline{toc}{chapter}{Agradecimientos}%

Esta tesis es el resultado de un camino recorrido con el apoyo de muchas personas a quienes debo mi más sincero agradecimiento.

En primer lugar, al profesor \textbf{Gustavo Osorio}, mi director de tesis, por su acompañamiento constante, su rigor académico y su confianza en este proyecto desde sus primeras ideas hasta su culminación. Su orientación fue determinante para mantener el rumbo y la calidad de la investigación.

A los estudiantes de pregrado y posgrado que hicieron parte activa del proceso de investigación, en especial a \textbf{Edward Grueso} y \textbf{Gustavo Pisso}, cuya colaboración en el laboratorio y sus aportes técnicos enriquecieron significativamente este trabajo.

A mi familia, por su apoyo incondicional, su paciencia durante las largas jornadas de estudio y su motivación permanente para alcanzar esta meta.

A mis colegas en la industria, quienes con su experiencia práctica, sus ideas y su visión del mundo real nutrieron la perspectiva aplicada de esta investigación. Sus conversaciones y retroalimentación fueron invaluables para conectar la academia con los desafíos reales del sector energético.

Finalmente, a la \textbf{Universidad Nacional de Colombia}, sede Manizales, y al \textbf{Departamento de Ingeniería Eléctrica, Electrónica y Computación}, por brindarme el espacio, los recursos y el entorno académico necesarios para llevar a cabo esta investigación.
